% year: תשפ"ב
% type: מבחן
% semester: א
% moed: ב
% number: 3
% points: 17
% categories: taylor
% source: exams/מבחנים/תשפב/מועד ב תשפב.pdf

\begin{question}
(17 נק') העריכו בקירוב את הערך של $\sqrt[4]{e}$ עם שארית הקטנה מ-$\frac{1}{1000}$. עליכם לנמק היטב את תשובתכם.
\end{question}

\begin{hint}
$\sqrt[4]{e}=e^{1/4}$. השתמשו בפולינום מקלורן של $e^x$ ובשארית לגרנז', ומצאו את הסדר $n$ הקטן שמבטיח שגיאה קטנה מ-$\frac{1}{1000}$.
\end{hint}

\begin{hint}
כדי לחסום את $e^c$ עבור $0<c<\frac14$ אפשר להשתמש ב-$e^c<e<3$.
\end{hint}

\begin{solution}
נכתוב $\sqrt[4]{e}=e^{1/4}$ ונשתמש בפולינום מקלורן של $f(x)=e^x$. כל הנגזרות הן $f^{(k)}(x)=e^x$, ולכן $f^{(k)}(0)=1$ ו-
\[ P_n(x)=\sum_{k=0}^n\frac{x^k}{k!}. \]
לפי משפט טיילור עם שארית לגרנז', לכל $n$ קיימת $c\in(0,\tfrac14)$ כך ש-
\[ R_n\left(\tfrac14\right)=\frac{e^c}{(n+1)!}\left(\tfrac14\right)^{n+1}. \]
כיוון ש-$e^x$ עולה, $e^c<e^{1/4}<e<3$, ולכן
\[ 0<R_n\left(\tfrac14\right)<\frac{3}{(n+1)!\,4^{n+1}}. \]
\begin{itemize}
  \item $n=2$: החסם הוא $\frac{3}{6\cdot64}=\frac{1}{128}$ — לא מספיק.
  \item $n=3$: החסם הוא $\frac{3}{24\cdot256}=\frac{1}{2048}<\frac{1}{1000}$ — מספיק.
\end{itemize}
(ואכן $n=2$ באמת אינו מספיק, ולא רק החסם שלו: $R_2=\frac{e^c}{6\cdot64}>\frac{1}{384}>\frac1{1000}$, כי $e^c>1$.)

לכן נבחר $n=3$:
\[ \sqrt[4]{e}\approx P_3\left(\tfrac14\right)=1+\frac14+\frac{1}{2\cdot16}+\frac{1}{6\cdot64}=1+\frac14+\frac1{32}+\frac1{384}=\frac{493}{384}\approx1.28385, \]
והשגיאה מקיימת $0<\sqrt[4]{e}-\frac{493}{384}<\frac{1}{2048}<\frac{1}{1000}$.

\textbf{תשובה:} $\boxed{\sqrt[4]{e}\approx\frac{493}{384}\approx1.2839}$ (הערך האמיתי $\approx1.28403$).
\end{solution}
